\begin{table*}[!t]
\caption{The ten most cited publications and their analytical focus.}
\label{tab:cited}
\centering
\footnotesize
\begin{tabularx}{\textwidth}{@{}L{3.3cm}YL{1.0cm}rrr@{}}
\toprule
First author, year & Signal and task & Type & TC & \shortstack{TC per\\year} & \shortstack{Normalised\\TC} \\
\midrule
Warrick, 2010 \cite{warrick2010} & Intrapartum CTG, hypoxia classification & Article & 89 & 5.24 & 2.78 \\
Ferrario, 2009 \cite{ferrario2009} & FHR variability, growth restriction & Article & 83 & 4.61 & 1.66 \\
Fergus, 2016 \cite{fergus2016} & EHG, preterm birth prediction & Article & 71 & 6.45 & 2.47 \\
Spilka, 2017 \cite{spilka2017} & Intrapartum FHR, outcome classification & Article & 68 & 6.80 & 2.51 \\
Fergus, 2018 \cite{fergus2018} & CTG, delivery mode classification & Article & 68 & 7.56 & 3.19 \\
Cömert, 2018 \cite{comert2018} & CTG, hypoxia assessment & Article & 61 & 6.78 & 2.86 \\
Cömert, 2019 \cite{comert2019} & CTG, hypoxia detection & Proc. & 60 & 7.50 & 2.48 \\
Ponsiglione, 2021 \cite{ponsiglione2021} & FHR processing and analysis & Review & 55 & 9.17 & 4.04 \\
Georgoulas, 2017 \cite{georgoulas2017} & FHR, evaluation by pH & Article & 55 & 5.50 & 2.03 \\
Jaba Deva Krupa, 2022 \cite{krupa2022} & Abdominal fECG, fetal QRS detection & Article & 48 & 9.60 & 3.07 \\
\bottomrule
\end{tabularx}
\abbrwide{Proc., proceedings paper. TC, total citations. Normalised TC uses the corpus mean for the publication year. CTG, cardiotocography. FHR, fetal heart rate. EHG, electrohysterography. fECG, fetal electrocardiography.}
\end{table*}
